% DATA PROVENANCE (generated 2026-11-02):
% - Channel constants read from the pinned NeuronBench package
%   (murphyk/neuronbench @ c354622458c460b419cab821d482c879f0578377) via
%   scripts/neuronbench_fully_observable.py:world_spec(), and cross-checked
%   against runs/neuronbench_fully_observable/data/manifest.json.
% - Train/held-out split: splits/neuron_first1.txt (z_rebound) for
%   meta-training; the other five worlds are held out (as in run 708907).
% Required packages in the paper preamble: booktabs, tabularx.

\begin{table}[t]
  \caption{Fully observable NeuronBench tasks. Each task is a noiseless
  regression from the external current $I_{\rm ext}$, membrane voltage $V$,
  and every channel's open fraction $\phi_c$ to the membrane vector field
  $\dot V$. All six worlds share the Hodgkin--Huxley sodium, potassium, and
  leak currents,
  $\dot V_{\rm HH} = I_{\rm ext} - 120\,\phi_{\rm Na}(V-50)
      - 36\,\phi_{\rm K}(V+77) - 0.3\,(V+54.4)$,
  with $\phi_{\rm Na}=m_{\rm Na}^3h_{\rm Na}$ and $\phi_{\rm K}=n_{\rm K}^4$,
  and differ in one added channel (or, for \texttt{na\_fatigue}, a modified
  $\phi_{\rm Na}$). Units are standard Hodgkin--Huxley (mV, ms,
  $\mu$A\,cm$^{-2}$, mS\,cm$^{-2}$, $C_m=1\,\mu$F\,cm$^{-2}$). Inputs are
  sampled by scrambled Sobol sequences over $I_{\rm ext}\in[-40,20]$,
  $V\in[-95,60]$, and $\phi_c\in[0,1]$. \texttt{z\_rebound} is used for
  meta-training; the other five worlds are held out.}
  \label{tab:neuronbench-tasks}
  \centering
  \small
  \setlength{\tabcolsep}{4pt}
  \renewcommand{\arraystretch}{1.15}
  \begin{tabularx}{\linewidth}{@{}l>{$}l<{$}>{\raggedright\arraybackslash}X@{}}
    \toprule
    Task & \multicolumn{1}{l}{Ground-truth equation} & Description \\
    \midrule
    \texttt{z\_rebound} & \dot V = \dot V_{\rm HH} - 4\,\phi_{\rm Z}(V-120)
      & Unidentified low-threshold inward current Z,
        $\phi_{\rm Z}=m_{\rm Z}^2h_{\rm Z}$; de-inactivated by
        hyperpolarization, producing depolarization block. \\
    \texttt{h\_sag} & \dot V = \dot V_{\rm HH} - 5\,\phi_{\rm H}(V+30)
      & Hyperpolarization-activated, non-inactivating inward current $I_h$,
        $\phi_{\rm H}=m_{\rm H}$; produces voltage sag and post-inhibitory
        rebound. \\
    \texttt{na\_fatigue} & \dot V = \dot V_{\rm HH}
      & No added current; $\phi_{\rm Na}=m_{\rm Na}^3h_{\rm Na}s_{\rm Na}$
        includes a slow inactivation gate, producing use-dependent
        spike-count run-down. \\
    \texttt{ca\_rebound} & \dot V = \dot V_{\rm HH} - 3.2\,\phi_{\rm T}(V-120)
      & Fast low-threshold T-type Ca$^{2+}$ current,
        $\phi_{\rm T}=m_{\rm T}^2h_{\rm T}$; de-inactivated by
        hyperpolarization, producing a rebound burst on release. \\
    \texttt{d\_type} & \dot V = \dot V_{\rm HH} - 9\,\phi_{\rm D}(V+77)
      & D-type K$^{+}$ current, $\phi_{\rm D}=m_{\rm D}h_{\rm D}$;
        de-inactivated by hyperpolarization, producing delayed or
        suppressed firing. \\
    \texttt{textbook\_M} & \dot V = \dot V_{\rm HH} - 2.5\,\phi_{\rm M}(V+77)
      & Textbook slow, non-inactivating M-type K$^{+}$ current,
        $\phi_{\rm M}=m_{\rm M}$; produces spike-frequency adaptation. \\
    \bottomrule
  \end{tabularx}
\end{table}
